%% ============================================================
%%  SOLUCIÓN — Ejercicio 2.2: Símbolos matemáticos
%%  Clase 1 — Después de diapositiva "Símbolos esenciales"
%%  Compilar: F5 (pdflatex)
%% ============================================================

\documentclass[12pt, a4paper]{article}
\usepackage[utf8]{inputenc}
\usepackage[T1]{fontenc}
\usepackage[spanish]{babel}
\usepackage{amsmath, amssymb}

\title{Símbolos matemáticos en economía}
\author{Tu Nombre}
\date{\today}

\begin{document}
\maketitle

\section{Expresiones econométricas}

Aquí están las expresiones del ejercicio usando LaTeX:

\subsection{Estimador OLS en forma matricial}
\[
  \hat{\beta}_{OLS} = (X'X)^{-1}X'Y
\]

\subsection{Supuesto de exogeneidad}
\[
  \mathbb{E}[\varepsilon_i | X_i] = 0
\]

\subsection{Varianza del estimador}
\[
  \text{Var}(\hat{\beta}) = \sigma^2(X'X)^{-1}
\]

\subsection{Condición de primer orden para máxima verosimilitud}
\[
  \frac{\partial \ln L(\theta)}{\partial \theta} = 0
\]

\subsection{Suma doble de cuadrados}
\[
  \sum_{t=1}^{T} \sum_{i=1}^{N} y_{it}^2
\]

\subsection{Fórmula del coeficiente de correlación}
\[
  \hat{\rho} = \frac{\sum_i (X_i - \bar{X})(Y_i - \bar{Y})}
                     {\sqrt{\sum_i(X_i-\bar{X})^2 \cdot \sum_i(Y_i-\bar{Y})^2}}
\]

\end{document}
